% =====================================================================
% Supplementary material -- adaptive-QMIX manuscript (Paper 1)
% RESULT-INDEPENDENT DRAFT. Same placeholder conventions as
% adaptive_qmix_main.tex. Source of truth: frozen commit
% 8c14f32988e92ddc6b1421d70258f5fbd2f2cc4d and Protocol Amendments 001
% and 002 (revised). The historical DCHF supplement (supplementary.tex)
% is preserved unchanged.
% =====================================================================
\documentclass[12pt]{article}

\usepackage[utf8]{inputenc}
\usepackage[T1]{fontenc}
\usepackage{lmodern}
\usepackage[margin=2.5cm]{geometry}
\usepackage{amsmath,amssymb,mathtools}
\usepackage{booktabs}
\usepackage{array}
\usepackage{tabularx}
\usepackage{xcolor}
\usepackage[inline]{enumitem}
\usepackage[colorlinks=true,linkcolor=blue,citecolor=blue,urlcolor=blue]{hyperref}

\newcolumntype{Y}{>{\raggedright\arraybackslash}X}
\newcolumntype{L}[1]{>{\raggedright\arraybackslash}p{#1}}
\renewcommand{\arraystretch}{1.15}
\newcommand{\TBD}[1]{\textcolor{red}{\texttt{\detokenize{#1}}}}
\newcommand{\Jp}{J_{\mathrm{primary}}}
\newcommand{\suppsec}[1]{\section*{#1}\addcontentsline{toc}{section}{#1}}

\title{Supplementary Material for\\
\emph{Adaptive Cooperative Traffic Signal Control with QMIX under
Short-Interval Green-Duration Decisions}}
\author{}
\date{}

\begin{document}
\maketitle

\noindent This supplement gives the exact specifications behind the
main paper: the executor and its timing identities (S-1), the network
and lane sets (S-2), traffic generation and random-number derivation
(S-3), vehicle accounting and the reward (S-4), the learning procedure
(S-5), the pretimed baseline search and its amendments (S-6), run
validity (S-7), the statistical procedure (S-8), behavioural metrics
(S-9), coordination diagnostics (S-10), provenance and chronology (S-11),
and result tables to be completed (S-12).

% ---------------------------------------------------------------------
\suppsec{S-1\quad Executor specification and timing identities}
% ---------------------------------------------------------------------

Both intersections share one decision clock with $\Delta t=5$~s. Let
$\phi_i\in\{H,V\}$ be the logical phase of intersection $i$ and $g_i$ its
elapsed green in seconds at a decision instant. Initially $\phi_i=H$ and
$g_i=0$ at $t=0$, and the first decision is taken at $t=0$. For each
decision:
\begin{itemize}
\item \textsc{extend}: the green of $\phi_i$ is shown for 5~s and
$g_i\leftarrow g_i+5$;
\item \textsc{switch}: the yellow of $\phi_i$ is shown for 3~s, then the
green of $\bar\phi_i$ for 2~s; afterwards $\phi_i\leftarrow\bar\phi_i$
and $g_i\leftarrow 2$.
\end{itemize}
Consequently a green entered by switching and extended $k\geq0$ times
lasts $2+5k$~s; the initial green lasts $5k$~s, $k\geq0$; a complete
H--V sequence of two greens entered by switching lasts $10+5(k_H+k_V)$~s;
and the time between phase changes at J1 and J2 is a multiple of 5~s.
Every change of phase is preceded by exactly 3~s of yellow, with no
all-red. The specification forbids a cycle length, an offset, a minimum
green and a maximum green as parameters of the primary controller. A
minimum-green wrapper would be a separate specification and is not
implemented. The specification is hashed
(\texttt{primary-adaptive-semantics-1.0}; see the provenance table of the
main paper).

% ---------------------------------------------------------------------
\suppsec{S-2\quad Network and lane sets}
% ---------------------------------------------------------------------

The nominal junction spacing is 300~m. The simulated lane lengths are
279.2~m for the central link in each direction, 289.6~m for the outer
main-corridor links and 189.6~m for all side links; the implementation
checks each length against these values at start-up (tolerance
0.001~m). Table~\ref{tab:lanes} lists the lanes used by the observation
and by the pressure controllers.

\begin{table}[!htbp]
\centering
\caption{Lane sets per intersection (SUMO lane identifiers).}
\label{tab:lanes}
\small
\begin{tabularx}{\textwidth}{@{}lL{3.4cm}Y@{}}
\toprule
 & J1 & J2 \\
\midrule
H incoming & \texttt{E0\_0/1}, \texttt{-E1\_0/1} & \texttt{E1\_0/1}, \texttt{E2\_0/1} \\
V incoming & \texttt{E4\_0/1}, \texttt{E5\_0/1} & \texttt{E6\_0/1}, \texttt{E7\_0/1} \\
H outgoing & \texttt{E1\_0/1}, \texttt{-E0\_0/1} & \texttt{-E1\_0/1}, \texttt{-E2\_0/1} \\
V outgoing & \texttt{-E4\_0/1}, \texttt{-E5\_0/1} & \texttt{-E6\_0/1}, \texttt{-E7\_0/1} \\
\bottomrule
\end{tabularx}
\end{table}

Max Pressure and P5 use the straight-through lane pairs (incoming lane,
receiving lane) of each movement, written out explicitly and checked
against the network connectivity at start-up. Their counts are the
number of vehicles on each lane, not the halting count, occupancy, the
QMIX observation or the reward.

% ---------------------------------------------------------------------
\suppsec{S-3\quad Traffic generation and random-number derivation}
% ---------------------------------------------------------------------

For each of the six streams (WE 900, EW 700, and four side streams of
300 vehicles), the departure seconds are drawn independently and
uniformly from $\{0,1,\dots,3599\}$ by a PCG64 generator seeded with a
\texttt{SeedSequence} over the experiment entropy, the
traffic-manifest namespace, the family code, the seed, the manifest
index and the stream rank. The family codes are: training 10,
development 20, benchmark design 30, benchmark validation 40, learner
validation 50, final test 60. Records are ordered by departure second,
then stream, then rank within stream, and every manifest has exactly
2{,}800 unique vehicles. The per-episode SUMO seed is derived from the
same identifiers in a separate namespace and folded into SUMO's signed
32-bit range. Learner randomness uses separate namespaces for the
initialization of each local network, the QMIX mixer, the two
exploration streams and replay sampling. Each exploration stream
consumes two draws per agent and decision whether or not the random
action is used, so random-number consumption is identical across
methods.

% ---------------------------------------------------------------------
\suppsec{S-4\quad Vehicle accounting and reward}
% ---------------------------------------------------------------------

A ledger tracks every scheduled vehicle as future, due-but-pending,
active, completed or exceptional (for example a teleport or an
unexpected disappearance). The network is \emph{cleared} when every
scheduled vehicle is completed and SUMO reports no remaining expected
vehicles. The reward of Eq.~(\ref{eq:reward_s}) is accumulated second by
second within each 5~s block:
\begin{equation}
r_k=-\sum_{\tau=5k+1}^{5k+5}\Bigl[\#\{v\ \text{active}: \mathrm{speed}_v(\tau)\leq0.1\}
+\#\{v: \text{scheduled departure}\leq\tau,\ \text{not inserted}\}\Bigr].
\label{eq:reward_s}
\end{equation}

% ---------------------------------------------------------------------
\suppsec{S-5\quad Learning procedure}
% ---------------------------------------------------------------------

\noindent\textbf{Procedure S-5 (official training run; identical for
QMIX, VDN and IDQN).}
\begin{enumerate}[leftmargin=1.2cm,label=\arabic*.]
\item Initialize the local networks from seed-derived streams, which are
identical across methods; QMIX also initializes its mixer. Copy the
online networks to the target networks; set $t\gets0$, $e\gets0$.
\item While $t<360{,}000$: load the training manifest with index $e$,
reset the simulation and observe $(o_1,o_2)$ and $s$. Then repeat:
\begin{enumerate}[label=(\alph*)]
\item $\varepsilon\gets1.0-0.95\cdot\min(t,250{,}000)/250{,}000$;
\item each agent takes its greedy local action, replaced with
probability $\varepsilon$ by a random action from its own stream;
\item execute 5~s; receive $r$, $(o'_1,o'_2)$, $s'$ and the flags
terminated, timeout and budget; store the joint transition; $t\gets t+1$;
\item if $t>5{,}000$: sample 256 transitions and perform one Double-DQN
update; hard-copy the targets every 500 updates;
\item if $t$ is a preregistered checkpoint, save the checkpoint;
\end{enumerate}
until the network clears, the 7{,}200~s cap is reached, or
$t=360{,}000$. Set $e\gets e+1$.
\end{enumerate}

The bootstrap mask uses the terminated flag only. Timeout and budget
truncation are stored as non-terminal, and the discount is
$\gamma^{\Delta t/5}$ with $\Delta t=5$~s. A transition without a valid
next observation is never stored. At the end of a full-budget run the
number of learner updates must equal 355{,}000, or the run is rejected.

% ---------------------------------------------------------------------
\suppsec{S-6\quad Pretimed baseline search and amendments}
% ---------------------------------------------------------------------

\paragraph{Plan construction} For cycle $C$ and split $f_H$ (in
thousandths), $g_H=\mathrm{round\_half\_up}\bigl(f_H(C-6)\bigr)$ computed in
exact decimal arithmetic, and $g_V=C-6-g_H$. A plan is admissible for
the classical search if $g_H\geq5$ and $g_V\geq5$. This rule applies to
the classical search only.

\paragraph{Coarse grids} Original: $C\in\{40,\dots,140\}$ in 10~s steps,
$f_H\in\{0.30,\dots,0.75\}$ in 0.05 steps, $\Delta\in\{0,5,\dots\}$ with
$\Delta<C$ (1{,}980 labels, all admissible). Lower-cycle audit (Amendment
001): $C\in\{20,25,30,35\}$ with the same splits and offset rule (208
admissible labels). Structural-floor audit (Amendment 002):
$C\in\{16,17,18,19\}$ with the same splits and offset rule, which gives
$\Delta\in\{0,5,10,15\}$ for each of these cycles
(Table~\ref{tab:floor}).

\begin{table}[!htbp]
\centering
\caption{Structural-floor audit grid (structural counts; no results).}
\label{tab:floor}
\small
\begin{tabular}{@{}cccll@{}}
\toprule
$C$ (s) & offsets $\Delta$ (s) & admissible labels & realized plans & $(g_H,g_V)$ (s) \\
\midrule
16 & 0, 5, 10, 15 & 8 & 4 & (5,5) \\
17 & 0, 5, 10, 15 & 12 & 8 & (5,6), (6,5) \\
18 & 0, 5, 10, 15 & 20 & 12 & (5,7), (6,6), (7,5) \\
19 & 0, 5, 10, 15 & 28 & 16 & (5,8), (6,7), (7,6), (8,5) \\
\midrule
total & & 68 & 40 & 200 design runs \\
\bottomrule
\end{tabular}
\end{table}

\paragraph{Realized-plan identity} The executor receives only
$(C,g_H,g_V,\Delta\bmod C)$. Labels with the same realized key are one
candidate. Every original label at $C\geq40$~s is a distinct plan
(1{,}980 coarse and 621 original fine labels, 2{,}547 in their union),
so realized-plan identity changes nothing in the original search. The
208 lower-cycle labels are 195 plans, and the 68 floor labels are 40
plans.

\paragraph{Fine neighbourhood} Around a centre $(C_0,f_0,\Delta_0)$:
$C\in\{C_0-10,\dots,C_0+10\}$ in 5~s steps within $[16,140]$ (originally
$[40,140]$); $f\in\{f_0-0.05,\dots,f_0+0.05\}$ in 0.025 steps within
$[0.25,0.80]$; $\Delta\in\{\Delta_0-10,\dots,\Delta_0+10\}$ in 1~s steps,
reduced modulo the candidate cycle; admissibility as above. Under
Amendment 002 the neighbourhood is generated around every label of the
combined coarse grid (1{,}980 + 208 + 68 labels) that realizes a winning
plan. The union is then deduplicated by realized plan. Around the five
best realized coarse plans this gives 1{,}359 labels and 792 realized
plans, of which 87 already had design runs and 705 are new (3{,}525 new
design runs). The ten best proceed to validation (50 runs).

\paragraph{Canonical naming and tie-breaks} A realized plan is named by
$\bigl(C,s^\ast,\Delta\bigr)$, where $s^\ast$ is the smallest integer
split in thousandths with $\mathrm{round\_half\_up}\bigl(s^\ast(C-6)/1000\bigr)=g_H$.
For fixed $C$, $g_H$ is non-decreasing in the split, so $s^\ast$ is
strictly increasing in $g_H$. Since $g_V=C-6-g_H$, lexicographic order of
canonical names equals lexicographic order of $(C,g_H,g_V,\Delta)$. The
frozen ranking tie-break (mean design $\Jp$, then name) and selection
tie-break (mean validation $\Jp$, time loss, cycle, then name) therefore
depend only on the realized plan. The driver re-sorts on the realized
key at every ranking and stops if the orders differ.

\paragraph{Replication gate (Amendment 001)} Five original candidates
were re-run through the lower-cycle audit path on the five design seeds.
Validity had to be identical, and mean $\Jp$ and mean time loss had to
agree within $10^{-12}$; all five candidates passed.

\paragraph{Evidence reuse} Each evaluated label has exactly one evidence
root, fixed by the grid that defined it. A run is reused only if its
sealed completion record verifies against the current manifests and the
full environment identity (source commit, configuration and network
hashes, campaign and log-schema versions), and if its run manifest
records the frozen commit, a clean working tree and one common software
stack. Labels of one realized plan that both have runs must agree per
seed within $10^{-12}$. A new run whose episode ends in
\textsc{clearance\_failure} receives a sealed terminal record and is not
re-simulated.

% ---------------------------------------------------------------------
\suppsec{S-7\quad Run validity gates}
% ---------------------------------------------------------------------

A run is admissible only if: all 2{,}800 scheduled vehicles are
accounted for; the status is cleared; there is no unexpected
disappearance; there is no unresolved teleport or accounting error; the
one-second evaluation log is complete; scientific analysis is permitted
for the run kind; and $\Jp$ and the required secondary outcomes (time
loss, travel time, waiting time, departure delay, stops) are finite. A
failed run contributes NaN and is never imputed. The statistical code
refuses non-finite inputs.

% ---------------------------------------------------------------------
\suppsec{S-8\quad Statistical procedure}
% ---------------------------------------------------------------------

\noindent\textbf{Procedure S-8 (paired crossed two-factor bootstrap for
one comparison).}
\begin{enumerate}[leftmargin=1.2cm,label=\arabic*.]
\item Form $D\in\mathbb R^{10\times10}$ with rows for training seeds
101--110 and columns for final seeds 3001--3010.
\item Seed a fixed generator from the protocol namespace and the matrix
shape.
\item For $b=1,\dots,10{,}000$: draw 10 row indices $R$ and,
independently, 10 column indices $K$ with replacement, and set
$m_b=\frac{1}{100}\sum_{r\in R}\sum_{c\in K}D_{rc}$, keeping
multiplicities (Cartesian product).
\item The 95\% interval is the 2.5th--97.5th percentile range of
$\{m_b\}$; the decision uses ``upper limit $<0$''.
\end{enumerate}

Additional quantities: $\bar D_{i\cdot}$ (per training seed);
$d_z=\overline{\bar D_{i\cdot}}/\mathrm{sd}(\bar D_{i\cdot})$ with the
$n-1$ denominator; the fraction of $i$ with $\bar D_{i\cdot}<0$; and the
training-success count, meaning runs completed under contract with a
validly selected checkpoint, with a 95\% Wilson score interval.
% TODO-LITERATURE: bootstrap, crossed-design resampling, d_z and Wilson
% references (see main paper).

% ---------------------------------------------------------------------
\suppsec{S-9\quad Behavioural metrics}
% ---------------------------------------------------------------------

A row of the one-second phase log stamped $t$ describes $[t-1,t)$. A
\emph{green spell} of movement $m$ is a maximal run of seconds with
actual phase $m$. The spell in progress at $t=0$ counts as a service
spell. A \emph{non-green spell} (service deprivation) is a maximal run of
seconds without green for $m$, including the conflicting green and the
yellow; spells still open at the end of the window are marked
right-censored. The \emph{emergent cycle} is the interval between
consecutive H-green starts at the same intersection. Reported per
intersection and movement over the demand window $[0,3600)$~s, with the
full episode as secondary: green-spell distribution (including the count
and share of 2~s greens); non-green spell mean, median, 95th percentile
and maximum; services per hour; green share among green time; yellow
fraction of elapsed time; queue 95th percentile and maximum; switch
count, rate and fraction of decisions (not applicable to pretimed
plans); and emergent-cycle mean, median, 95th percentile and standard
deviation. No numeric starvation or flickering threshold is defined
(see the behavioural-review subsection of the main paper).

% ---------------------------------------------------------------------
\suppsec{S-10\quad Coordination diagnostics}
% ---------------------------------------------------------------------

Let $g_i(t)\in\{0,1\}$ indicate H green at intersection $i$.
\begin{itemize}
\item Phase cross-correlation:
$\rho(\ell)=\mathrm{corr}\bigl(g_1(t),g_2(t+\ell)\bigr)$ for
$|\ell|\leq120$~s; a positive $\ell$ means that J2 follows J1.
\item Green-start lag: for each H-green start at J1, the time to the
first H-green start at J2 at or after it (and symmetrically for EW).
\item GAD50: for each direction on the central link, the fraction of
vehicle crossings of a virtual detector 50~m upstream of the downstream
stop line (at 229.2~m along the 279.2~m lane) that occur during
downstream H green.
\item Projected arrival on green: detector crossings projected to the
stop line at the free-flow speed of 13.89~m/s. It is a free-flow proxy,
never an observed arrival-on-green measure.
\item Downstream stops of vehicles that crossed the detector.
\end{itemize}
Rolling versions use 300~s windows advanced every 60~s. Because the
learned controllers act on a common 5~s grid, lags are interpreted at
5~s resolution. \TBD{TBD_AFTER_FINAL_EVALUATION}.

% ---------------------------------------------------------------------
\suppsec{S-11\quad Provenance and chronology}
% ---------------------------------------------------------------------

\begin{table}[!htbp]
\centering
\caption{Protocol chronology (to be completed with dates and hashes).}
\label{tab:chronology}
\scriptsize
\begin{tabularx}{\textwidth}{@{}L{4.6cm}Y@{}}
\toprule
Event & Record \\
\midrule
Frozen implementation & commit \texttt{8c14f32988e92ddc6b1421d70258f5fbd2f2cc4d} \\
Original classical design, validation and freeze & fixed-offset and fixed-timing tracks; fixed-timing winners on the $C=40$~s bound \\
QMIX seed 101 (pre-amendment) & launched before Amendment 001 was documented; no checkpoint selected; excluded from inference \\
Amendment 001 (lower-cycle audit) & SHA-256 \texttt{edb9ec98281b6a9588fae96cfdf982eb90984b830b7fa815ad058160b3ba020a}; replication gate passed; fixed-timing track reopened \\
Amendment 002, first draft & withdrawn before execution (label-based deduplication); archived \\
Amendment 002 (revised), sealed & sealed before the structural-floor and reopened fine-search simulations; lower-cycle evidence reused only after provenance and realized-plan equivalence checks; hashes \TBD{TBD_AFTER_FINALISE_COARSE} \\
Structural-floor audit and combined coarse ranking & completed; 40 floor plans on design seeds; five best realized plans in Table~\ref{tab:coarse_design} \\
Reopened fine search and timing validation & 705 new realized plans (3{,}525 design runs), then 50 validation runs; \TBD{TBD_AFTER_TIMING_VALIDATION} \\
Amended baseline freeze & \TBD{TBD_AFTER_TIMING_VALIDATION} \\
Official training (30 runs) & \TBD{TBD_AFTER_LEARNER_TRAINING} \\
Seed-101 checkpoint-hash comparison & \TBD{TBD_AFTER_LEARNER_TRAINING} \\
Checkpoint freeze and behavioural reviews & \TBD{TBD_AFTER_LEARNER_VALIDATION} \\
Pre-final audit, final-test unlock, final evaluation & \TBD{TBD_AFTER_FINAL_EVALUATION} \\
\bottomrule
\end{tabularx}
\end{table}

% ---------------------------------------------------------------------
\suppsec{S-12\quad Result tables to be completed}
% ---------------------------------------------------------------------

\begin{table}[!htbp]
\centering
\caption{Combined coarse ranking, five best realized plans (design seeds
2001--2005). The design-set $\Jp$ values are selection-optimistic: these
plans were retained because they scored best on these seeds, so the
values are not estimates of baseline performance and must not be
compared with held-out results.}
\label{tab:coarse_design}
\small
\begin{tabular}{@{}clcl@{}}
\toprule
Rank & Realized $(C,g_H,g_V,\Delta)$ & design $\Jp$ (s) & Active constraints \\
\midrule
1 & (18, 7, 5, 10) & 2.809142857143 & $g_V=5$ \\
2 & (19, 8, 5, 10) & 2.825785714286 & $g_V=5$ \\
3 & (20, 9, 5, 10) & 3.042214285714 & $g_V=5$ \\
4 & (25, 12, 7, 0) & 3.280071428571 & none \\
5 & (25, 13, 6, 0) & 3.289928571429 & none \\
\bottomrule
\end{tabular}
\end{table}

\begin{itemize}
\item Fine-search counts per winner and validation table of the ten
finalists:\newline \TBD{TBD_AFTER_TIMING_VALIDATION}.
\item Checkpoint-validation matrices ($14\times5$ per method and seed) and
selected checkpoints: \TBD{TBD_AFTER_LEARNER_VALIDATION}.
\item Per-seed final-test $\Jp$ matrices, secondary outcomes, behavioural
distributions and coordination diagnostics:
\TBD{TBD_AFTER_FINAL_EVALUATION}.
\item High-demand slice: \TBD{OPTIONAL_ROBUSTNESS_AFTER_PRIMARY_GATE}.
\end{itemize}

\end{document}
